\textbf{Actividades de difusión y extensión}

\begin{enumerate}

\item \textbf{"Caravana del Cerebro"}, 25 de abril, Tipo: Semana del cerebro, 2024.

\item \textbf{"Neuro-Emociónate"}, CAC-UNAM, del 19 al 23: , Tipo: Semana del Cerebro, marzo de 2024.

\item  \textbf{"Semana del Cerebro"}, Centro Educativo y Cultural ‘Manuel Gómez Morin’, Querétaro, Qro., 19-23 marzo, Tipo: Semana del 
cerebro, marzo de 2019.

\item  \textbf{"¡Muévete ... Mantén activo tu Cerebro!"}, Semana del Cerebro, Tipo: En la semana del cerebro, marzo de 2016.

\item  \textbf{"Feria de las Ciencias y las Humanidades, El cerebro del buen comer"}, Tipo: Conferencia, octubre de 2013.

\item  \textbf{“¿Qué tengo en mi cabeza?"}, Participación de Carteles, Tipo: Semana del cerebro, marzo de 2012.

\item  \textbf{"Semana del Cerebro, CerebrAte"}: , Tipo: En la semana del cerebro, marzo de 2011.

\item  \textbf{"Juego de Rompecabezas del Cerebro"}, Participación en las actividades de la Semana del Cerebro, Marzo, 16-20, Tipo: En 
la 
semana del cerebro, 2010.

\item  \textbf{"Puedes ver los que pasa en tu cerebro...Sí!"}: La Sociedad Mexicana de Ciencias Fisiológicas y El Capítulo de la Ciudad 
de 
México de la "Society for Neuroscience",Marzo, 15-21, Tipo: Semana del Cerebro, 2010.

\item  \textbf{Organización de la "Celebración Internacional del Cerebro"}, Tipo: Semana del Cerebro, abril de 2006.

\item  \textbf{Participación en la "Celebración Internacional del Cerebro"}, Tipo: Semana del Cerebro, marzo de 2005.

\end{enumerate}
